\part*{Appendices}
\addcontentsline{toc}{part}{Appendices}

\section{Notation registry and reconciliation with Paper~A}\label{app:notation}

This appendix fixes every symbol used in the paper, and records how the notation differs from that of the accepted companion manuscript (Paper~A). Nothing in Paper~B depends on reading Paper~A; the reconciliation is given only so that a reader of both can translate.

\begin{center}\small
\begin{tabular}{@{}p{3.1cm}p{11.4cm}@{}}
\toprule
Symbol & Meaning\\\midrule
$n$ & number of coordinates (length of the code); the box is $[0,1]^n$.\\
$M$, $C=\ker M$ & binary matrix and its code $C\subseteq\F_2^n$; $k=\dim C$ is the \emph{nullity}.\\
$\Lambda_M$ & the lattice $\{\alpha\in\Z^n: M\alpha\equiv0\bmod2\}$.\\
$F=C\setminus\{0\}$ & the fibre: the integer points of the box in $\Lambda_M$ other than the origin.\\
$\supp w$, $\wt w$ & support and Hamming weight of $w$.\\
$\Circ(C)$ & the set of circuits (minimal supports of nonzero words); words and supports are identified.\\
$S$, $U(S)$, $\spn S$ & a family of pairwise disjoint circuits, its union, and the $2^{|S|}$ words obtained by taking unions of members.\\
$Z(S,x)$ & the cell with base point $x\in C\setminus\spn S$ and generators $g_j=(1-2x_i)_{i\in S_j}$.\\
$p$, $X$ & reduced length (number of series classes) and the reduced point set $X\subseteq\AG(k-1,2)$.\\
$\Part$ & a partition of $\one$ into circuits; $V_\Part=\spn\Part$, $m_\Part=\dim V_\Part$.\\
$n_m$, $n_2$ & number of partitions of $\one$ into exactly $m$ circuits; $n_2$ counts the two-block partitions.\\
$N$ & number of cells through the centre of the box (Proposition~\ref{prop:centre-count}).\\
$\rmin$ & minimum witness rank: the least $\dim\operatorname{aff}(Z\cup Z')$ over collisions $(Z,Z')$.\\
$E_p$ & the even-weight code of length $p$ (cut space of the cycle $C_p$).\\\bottomrule
\end{tabular}
\end{center}

\paragraph{Differences from Paper~A.}
(i) Paper~A indexes coordinates by a number of primes that it also calls $k$; here $n$ is the length and $k$ is always the nullity. (ii) Paper~A studies planar modular lattices $\Gamma_{d,s}$ arising from two-active-prime occurrence fibres; Paper~B works with the relation lattice $\Lambda_M$ of an arbitrary binary matrix and, for the minimal non-embedded boxes, the relevant parity vector is the all-ones word. (iii) Terminology: ``collision'' and ``witness rank'' are not formally introduced in Paper~A; Paper~B defines them precisely (Section~\ref{sec:cells}) and uses terminology consistent with the GEO-I non-embedding releases. $\rmin$ is the minimum over collisions and is not the ambient dimension, the nullity, or the rank of every witness. (iv) Terminology: ``arithmetic occurrence geometry'' is the defined term of Paper~A; here it is used only through the specialization of Section~\ref{sec:aog}. (v) All theorem numbers of this paper are new; Appendix~\ref{app:crosswalk} maps them to release identifiers.

\section{The Graver basis of \texorpdfstring{$\Lambda_M$}{Lambda\_M}: further detail}\label{app:B}

Proposition~\ref{prop:graver} states that the Graver basis of $\Lambda_M$ consists of the $\pm1$ sign vectors supported on circuits together with the vectors $\pm2e_i$ for coordinates $i$ that are not singleton circuits. We record the correspondence with the standard theory \cite{Graver1975,Sturmfels1996,Hemmecke2003} and the two points that are easy to misstate.

\paragraph{Conformal order.} For $u,v\in\Z^n$ write $u\conf v$ if $u_iv_i\ge0$ and $|u_i|\le|v_i|$ for all $i$. The Graver basis is the set of $\conf$-minimal nonzero elements of $\Lambda_M$. An element of the box fibre is a $0/1$ vector $w\in C$; its sign-extensions $\pm w$ are the lattice vectors it represents.

\paragraph{Why only circuits and doubles appear.} Let $\alpha\in\Lambda_M$ be $\conf$-minimal. If some $|\alpha_i|\ge3$ then $\alpha-2\sgn(\alpha_i)e_i\in\Lambda_M$ is conformally below $\alpha$, a contradiction; so $|\alpha_i|\le2$. If $|\alpha_i|=2$ then $\alpha\conf2e_i$ forces $\alpha=\pm2e_i$. Otherwise $\alpha\in\{0,\pm1\}^n$, its support is the support of a nonzero codeword, and minimality of $\alpha$ among sign-compatible vectors is exactly minimality of the support, i.e.\ the support is a circuit. Finally $\pm2e_i\in\Lambda_M$ is $\conf$-minimal unless $\pm e_i\in\Lambda_M$, which happens exactly when $\{i\}$ is a circuit.

\paragraph{Pitfall.} The doubles $\pm2e_i$ are Graver elements but are \emph{not} corners of any cell of the box: their endpoints leave $[0,1]^n$. They are irrelevant to the geometry of the box, and Proposition~\ref{prop:cells} uses only the circuit part. Statements that quantify over ``all Graver elements'' must therefore be restricted before being applied to the $0/1$ model.

\section{Cells, collisions and the centre count: formal definitions}\label{app:C}

\paragraph{Cells.} For a family $S=\{S_1,\dots,S_r\}$ of disjoint circuits and $x\in C\setminus\spn S$, the cell $Z(S,x)$ is the convex hull of the $2^r$ points $x+\sum_{j\in J}g_j$ with $J\subseteq[r]$ (with $g_j=(1-2x_i)_{i\in S_j}$ extended by zero). By Lemma~\ref{lem:cellpair} a cell is determined by the pair $(S,\,x+\spn S)$; the vertices are the words of the coset $x+\spn S$ (mod $2$ the coset is closed under addition of members of $\spn S$, and the $0/1$ representatives are exactly the vertices).

\paragraph{Collision.} Two distinct cells of positive dimension \emph{collide} if their relative interiors intersect (Lemma~\ref{lem:relint}). A box is \emph{embedded} iff no collision exists, equivalently (Proposition~\ref{prop:facetoface}) iff any two cells meet in a common face (possibly empty).

\paragraph{Centre count.} Every collision of a minimal non-embedded box is centred at $\tfrac12\one$ (Theorem~\ref{thm:centre} together with Theorem~\ref{thm:F1}). The cells whose relative interior contains the centre are in bijection with the pairs $(\Part,\,x+V_\Part)$ with $x\notin V_\Part$; counting cosets gives $2^{k-m_\Part}-1$ cells for each partition $\Part$. Summing yields $N=\sum_\Part(2^{k-m_\Part}-1)$, and since all $N$ cells contain the centre in their relative interior any two of them collide, so there are exactly $\binom N2$ collisions (Corollary~\ref{cor:collcount}). Example: $E_4$ has $k=3$ and three partitions of $\one$ into two circuits, $\{12,34\},\{13,24\},\{14,23\}$, each with $m=2$; each gives $2^{3-2}-1=1$ cell, so $N=3$ and there are $\binom32=3$ collisions.

\section{The circuit-partition criterion: variants}\label{app:D}

The proof of Theorem~\ref{thm:E1} rests on Lemma~\ref{lem:notunion}: an $\F_2$-vector space is not the union of two proper subspaces. Three comments on its use.

\paragraph{Where the lemma enters.} In (ii)$\Rightarrow$(i) of Theorem~\ref{thm:E1}, Lemma~\ref{lem:notunion} supplies a base point $x\notin\spn S\cup\spn T$; this is the only place where the field $\F_2$ is used, and it is why the two cells share a base point and hence satisfy the centre criterion.

\paragraph{The case $S=T$.} In (i)$\Rightarrow$(ii) the collision may a priori consist of two cells over the \emph{same} circuit family $S$ with different cosets. The proof shows this cannot happen without a second partition of $U$: the difference $c=x\oplus y\in C\setminus\spn S$ is supported in $U$, and decomposing $c$ and $\one_U\oplus c$ into circuits gives a partition $S'\ne S$. 

\paragraph{Collisions at a centre.} By Theorem~\ref{thm:centre}, collisions of a minimal box are exactly the pairs of distinct cells through the centre; their number is $\binom N2$ (Appendix~\ref{app:C}). The criterion therefore does more than detect non-embedding: it enumerates all collisions.

\section{Shortening, minimality and full support: proofs of the lemmas}\label{app:E}

Let $C_{(i)}=\{w|_{[n]\setminus\{i\}}: w\in C,\ w_i=0\}$ be the shortening of Definition~\ref{def:shortening} (equivalently, the kernel of $M$ with column $i$ deleted). Lemma~\ref{lem:shortcirc}(a) says that the circuits of $C_{(i)}$ are exactly the circuits of $C$ avoiding $i$. The proof is a direct comparison of supports: words of $C_{(i)}$ are the words of $C$ vanishing at $i$; any word of $C$ with support contained in a circuit avoiding $i$ also vanishes at $i$, so minimality transfers in both directions.

\paragraph{Descent (Proposition~\ref{prop:descent}).} If $C$ is non-embedded and some single shortening is again non-embedded, replace $C$ by it; the number of coordinates strictly decreases, so the process stops at a coordinatewise-minimal non-embedded box. Collisions in $C_{(i)}$ lift to collisions in $C$ with the same affine-span dimension, hence the minimum witness rank of $C$ is at most that of the minimal sub-box, and Part~III bounds the latter.

\paragraph{Full support (Theorem~\ref{thm:F1}).} Let $C$ be minimal and suppose a doubly-partitioned word $U$ misses a coordinate $i$. Then both partitions are circuit partitions of $U$ inside $C_{(i)}$ by Lemma~\ref{lem:shortcirc}, so $C_{(i)}$ is non-embedded by Theorem~\ref{thm:E1}, contradicting minimality. Hence $U=[n]$, and $\one=U\in C$.

\section{Affine reduction and the classification method}\label{app:F}

\paragraph{Generator matrices.} Choose a basis of $C$ and let $G$ be the $k\times n$ generator matrix; its columns are points of $\F_2^k$. Since $\one\in C$, there is a functional that takes the value $1$ on every column, so the columns lie in an affine hyperplane $H\cong\AG(k-1,2)$ (Theorem~\ref{thm:affine}). Identical columns form series classes and are merged (Proposition~\ref{prop:series}); the result is the set $X\subseteq H$ of $p$ distinct points, spanning $\F_2^k$.

\paragraph{Equivalence.} Two minimal boxes are equivalent if their reduced point sets differ by an element of $\AGL(k-1,2)$ (Lemma~\ref{lem:equivalence}). The group has order $2^{k-1}\prod_{i=0}^{k-2}(2^{k-1}-2^i)$, namely $24,\ 1344,\ 322\,560$ for $k=3,4,5$. For each $k\le5$ \texttt{scripts/classify.py} enumerates all $2^{2^{k-1}}$ subsets of $H$, selects those that are spanning with every shortening embedded and the full box non-embedded, and groups them into orbits using union--find under generators of $\AGL(k-1,2)$; the orbit sizes divide the group order and the Burnside count $\sum|\mathrm{Fix}|/|G|$ confirms the number of orbits. The result is Theorem~\ref{thm:class35}: $1$, $3$, $12$ classes for $k=3,4,5$.

\paragraph{Simplex normal form.} For $p\le8$ one may by an affine change of coordinates assume that $X$ contains the standard simplex $\{e_1,\dots,e_k\}$ (Lemma~\ref{lem:simplex}); the remaining $p-k$ points are odd-weight vectors of weight $\ge3$, and the finite check of Proposition~\ref{prop:small} runs over all such completions.

\section{Rigidity and the vector-space partition: a worked example}\label{app:G}

Take $k=4$, $p=6$ and the point set $X=\{1000,1001,1010,1011,1100,1101\}\subseteq\AG(3,2)$ (columns of the generator matrix, written in binary). The code has $2^4-1=15$ nonzero words and eleven circuits:
\[
12,\ 34,\ 56,\ 135,\ 136,\ 145,\ 146,\ 235,\ 236,\ 245,\ 246 .
\]
The partitions of $\one=123456$ into circuits are
\begin{align*}
&\{12,34,56\},\qquad \{135,246\},\ \{145,236\},\ \{136,245\},\ \{146,235\},
\end{align*}
i.e.\ $n_3=1$, $n_2=4$ (the profile $2^4,3^1$ of Table~\ref{tab:classlist}). The subspaces $V_\Part$ have dimensions $m=3$ for the three-block partition and $m=2$ for each of the four others, and they pairwise intersect in $\{0,\one\}$ (Theorem~\ref{thm:vsp}).

\paragraph{Counting identity.} $\sum_\Part(2^{m_\Part}-2)=4\cdot2+1\cdot6=14=2^4-2$.

\paragraph{Parseval identity.} $\sum_\Part(2^{m_\Part}S_\Part-2p^2)$ with $S_\Part=\sum_{c\in\Part}|c|^2$: the four two-block partitions have $S=9+9=18$ and contribute $4\cdot18-72=0$ each; the three-block partition has $S=4+4+4=12$ and contributes $8\cdot12-72=24$. Total $24=p(2^k-2p)=6\cdot4$. Both identities are thus checked by hand on this example; Table~\ref{tab:checks} records the exact verification for all classified classes.

\paragraph{Rigidity (Theorem~\ref{thm:rigid}).} Every word other than $\one$ has a unique circuit decomposition: for instance $1234=12\sqcup34$ is the only decomposition of this word, and $135$ is a circuit.

\section{Counting and Parseval: derivations}\label{app:H}

The derivations of Theorems~\ref{thm:counting} and~\ref{thm:parseval} are given in full in Part~III. We restate the chain of equalities in one place and list the exact identities verified on every classified class.

\begin{enumerate}
\item Partition $C\setminus\{0,\one\}$ into the sets $V_\Part\setminus\{0,\one\}$ (Theorem~\ref{thm:vsp}); counting elements gives $\sum_\Part(2^{m_\Part}-2)=2^k-2$.
\item Sum the squared weight over $C$: $\sum_{w\in C}\wt(w)^2=2^{k-2}p(p+1)$ for a reduced code with no repeated coordinate (Lemma~\ref{lem:moment}).
\item Expand the same sum block by block (Lemma~\ref{lem:block}): grouping the words of $C\setminus\{0,\one\}$ by the block $V_\Part$ containing them and subtracting the common contribution of $0$ and $\one$ yields $\sum_\Part(2^{m_\Part}S_\Part-2p^2)=p(2^k-2p)$.
\end{enumerate}

\begin{table}[ht]
\centering\small
\begin{tabular}{rrlcc}
\toprule
$k$ & $p$ & profile & counting & Parseval\\\midrule
\identitytablebody
\bottomrule
\end{tabular}
\caption{Exact values of the two sides of the counting and Parseval identities for the classes with $k\le5$ (output of \texttt{scripts/geo2.py::check\_identities}). The identities are also verified for the $215$ classes with $k=6$ (Table~\ref{tab:checks}).}
\end{table}

\section{Inequalities and the finite exceptional range}\label{app:I}

Part~III shows $n_2\ge p-1$ for $p\ge9$ from the threshold lemma $p(2^m/m-2)\ge2^m-2$ for $p\ge9$, $3\le m\le p/2$. The three smallest cases of the function $f(m)=m(2^m-2)/(2^m-2m)$ are $f(3)=9$, $f(4)=7$, $f(5)=75/11$ (Table~\ref{tab:threshold}), so only $m=3$ is binding and it is binding exactly at $p=9$.

\paragraph{The range $p\le8$.} For $p\le8$ the threshold inequality does not exclude minimal boxes whose partitions have large blocks, and the universality of W2 is established instead by exhaustive enumeration in simplex normal form (Proposition~\ref{prop:small}, Table~\ref{tab:smallp}). The counts of candidate completions and of minimal boxes among them are reproduced from \texttt{scripts/small\_p.py} and \texttt{logs/small\_p.json}. Minimal boxes with $p\le8$ occur only for $(k,p)\in\{(3,4),(4,6),(4,7),(4,8),(5,8)\}$ and each has $n_2\ge2$. The code $E_8$ ($k=7$, $p=8$) is not minimal (Figure~\ref{fig:E8}).

\section{Construction of rank-three witnesses}\label{app:J}

Let $\Part=\{b_1,b_2\}$ and $\Part'=\{c_1,c_2\}$ be two distinct partitions of $\one$ into two circuits (these exist in every minimal box, Theorem~\ref{thm:W2}). Choose $x\in C\setminus(V_\Part\cup V_{\Part'})$, which exists by Lemma~\ref{lem:notunion}. The cells $Z(\Part,x)$ and $Z(\Part',x)$ are parallelograms that collide (Theorem~\ref{thm:E1}). Their generators are the restrictions of the single vector $s=(1-2x_i)_{i\in[n]}$ to the four blocks, so $g_1+g_2=s=h_1+h_2$ is a linear relation among four vectors and the rank is at most $3$; Lemma~\ref{lem:LOW} gives at least $3$ (Theorem~\ref{thm:L5}). In the table $x$ is recorded as the set of coordinates carrying the value $1$ in a chosen base point, and the rank is computed independently by exact Fraction elimination, for the first listed minimal class of each $(k,p)$ with $k\le5$.

\begin{table}[ht]
\centering\scriptsize
\resizebox{\textwidth}{!}{\begin{tabular}{rrllcc}
\toprule
$k$ & $p$ & partition $\Part$ & partition $\Part'$ & base $x$ & rank\\\midrule
\witnesstablebody
\bottomrule
\end{tabular}}
\caption{Rank-three witnesses constructed as in Appendix~\ref{app:J} (first listed class per $(k,p)$).}
\end{table}

\subsection*{Partition-profile types}\label{app:J2}

For completeness we list the multiset of block sizes occurring over all classified classes. A \emph{type} is the multiset of $m_\Part$-values over partitions with $m\ge3$ blocks of circuits (partitions with two blocks have $m=2$ and appear in every profile).

\begin{table}[ht]
\centering\small
\begin{tabular}{llr}
\toprule
$k$ & type & number of classes\\\midrule
\typestablebody
\bottomrule
\end{tabular}
\caption{Types occurring among the classified minimal classes ($k\le6$), as computed by \texttt{scripts/gen\_tables.py} from the exhaustive and source-reported classifications. The rows with $k=6$ rely on the source-reported classification (certificates replayed with the independent engine).}
\end{table}

\section{Classification methods and certificates}\label{app:K}

\paragraph{$k\le5$: independent exhaustive enumeration.} \texttt{scripts/classify.py} enumerates all subsets of the affine space, tests minimality with the engine \texttt{geo2.py} (circuits by brute force from the definition, E1 criterion, single-shortening test), and groups subsets into orbits. The output (\texttt{logs/classify.json}) matches the release counts $1,3,12$.

\paragraph{$k=6$: source-reported.} The enumeration over $2^{32}$ subsets is not repeated here. The release GEO-I v010 (\texttt{BINARY6.md}) reports $215$ classes and supplies certificates \texttt{C6\_rank6.json} with the checker \texttt{checker10.py}. In this work the checker was run (output in \texttt{logs/source\_checkers.txt}) and the $215$ representatives were re-examined with the independent engine: each is spanning, non-embedded, with all shortenings embedded, and satisfies the counting and Parseval identities (\texttt{scripts/verify6.py}, \texttt{logs/verify6.json}). Pairwise inequivalence of the $215$ representatives was checked by invariants (\texttt{scripts/inv6*.py}) together with an equivalence search (\texttt{scripts/equiv6.py}). \emph{Completeness} of the list rests on the source and is not reproved here; nothing in the universal proof depends on it.

\paragraph{$p\le8$.} Proposition~\ref{prop:small} is a complete finite verification using the simplex normal form; its details are in \texttt{scripts/small\_p.py} and \texttt{logs/small\_p.json}.

\section{Reproducibility}\label{app:repro}

\paragraph{Pinned sources.} All citations to the GEO-I releases v001--v012 refer to the archived research-release repository at commit \geoicommit. The repository itself could not be fetched during the preparation of this manuscript; local copies of the release files supplied by the author were used, and the commit identifier was supplied by the author. These releases are archived research releases, not peer-reviewed publications.

\paragraph{File tree.}
\begin{center}\small
\begin{tabular}{@{}lp{10.5cm}@{}}
\texttt{tex/} & LaTeX sources (\texttt{main.tex}, \texttt{part1--5.tex}, \texttt{appendices.tex}, \texttt{refs.tex}, generated \texttt{gen\_data.tex}) and \texttt{figures/}.\\
\texttt{scripts/} & \texttt{geo2.py} (independent engine), \texttt{classify.py}, \texttt{crosscheck.py}, \texttt{lit\_checks.py}, \texttt{verify6.py}, \texttt{small\_p.py}, \texttt{inv6*.py}, \texttt{equiv6.py}, \texttt{gen\_tables.py}, \texttt{figs.py}.\\
\texttt{logs/} & JSON/text outputs of the above, including \texttt{source\_checkers.txt}.\\
\texttt{src/} & copies of the release certificates and checkers used (\texttt{C5\_rank5.json}, \texttt{C6\_rank6.json}, \texttt{C9\_partition.json}, \texttt{checker10.py}, \texttt{checker11.py}).\\
\end{tabular}
\end{center}

\paragraph{Commands.} With Python~3 and the packages \texttt{numpy} and \texttt{matplotlib}, run in \texttt{scripts/}: \texttt{python3 classify.py}, \texttt{python3 verify6.py}, \texttt{python3 small\_p.py}, \texttt{python3 gen\_tables.py}, \texttt{python3 figs.py}; then three passes of \texttt{pdflatex main.tex} in \texttt{tex/}. The engine uses only the Python standard library (exact rational elimination with \texttt{fractions.Fraction}).

\paragraph{Exact versus numerical.} All ranks and identities reported are exact (integer or rational arithmetic); no floating-point computation enters any claim. Floating point is used only in drawing the figures.

\paragraph{Environment.} Python 3.13, pdfTeX 3.141592653-2.6-1.40.25 (TeX Live 2023), Linux.

\section{Theorem crosswalk and provenance}\label{app:crosswalk}

``Release'' refers to the GEO-I archived release and file in which the statement is recorded; ``Source'' is the status reported there; ``Here'' says how the statement is established in this paper. ``New'' marks statements not found in the release files that were read.

\begingroup\scriptsize
\setlength{\tabcolsep}{3pt}
\begin{longtable}{@{}p{2.6cm}p{2.0cm}p{3.4cm}p{2.0cm}p{3.6cm}@{}}
\toprule
Statement here & Release id & Release / file & Source status & Status here\\\midrule\endhead
Prop.~\ref{prop:cells}, Thm.~\ref{thm:centre} & B1, B2 & v006 \texttt{OBSTRUCTION\_RANK.md} & proved & proved in text\\
Thm.~\ref{thm:E1} & G10-E1 & v010 \texttt{MATHEMATICS.md} \S1.1 & proved & proved in text\\
Prop.~\ref{prop:centre-count}, Cor.~\ref{cor:collcount} & -- & -- & -- & new; proved in text; enumerated for $k\le5$\\
Prop.~\ref{prop:descent} & G9-T5 & v009 \texttt{MATHEMATICS.md} \S2.4 & proved+checked & proved in text\\
Cor.~\ref{cor:nullity} & B-NULL$'$ & v008 \texttt{MINIMAL\_RANK.md} & proved & proved in text\\
Prop.~\ref{prop:series} & B-SERIES & v008 \texttt{MINIMAL\_RANK.md} & proved+checked & proved in text\\
Thm.~\ref{thm:F1} & G10-F1 (closes G8-C3) & v010 \texttt{MATHEMATICS.md} \S1.2 & proved & proved in text\\
Prop.~\ref{prop:nosingleton} & G9-Z, G11-L5c & v009 \S1; v011 \S1.1 & proved(+checked) & proved in text (one argument)\\
Thm.~\ref{thm:affine} & G10-A1 & v010 \texttt{MATHEMATICS.md} \S1.3 & proved & proved in text\\
Lem.~\ref{lem:LOW} & B-LOW (G8-T3a) & v008 \texttt{MINIMAL\_RANK.md} & proved & proved in text\\
Thm.~\ref{thm:L5} & G10-L5, B-COMMON & v010 \S1.5; v008 & proved & proved in text\\
Thm.~\ref{thm:C1p} & G10-C1$'$ & v010 \texttt{MATHEMATICS.md} \S1.5 & proved & proved in text\\
Thm.~\ref{thm:rigid} & G10-R & v010 \texttt{MATHEMATICS.md} \S1.4 & proved & proved in text\\
Prop.~\ref{prop:rigid-char} & -- & -- & -- & new; proved in text\\
Thm.~\ref{thm:vsp}, \ref{thm:counting} & G11-L5 & v011 \texttt{MATHEMATICS.md} \S1.2--1.3 & proved & proved in text\\
Thm.~\ref{thm:parseval} & G11-L5b & v011 \texttt{MATHEMATICS.md} \S1.4 & proved & proved in text\\
Thm.~\ref{thm:L6} & G11-L6 (sharp: G11-K1) & v011 \texttt{MATHEMATICS.md} \S1.5 & proved+checked & proved in text; equality case analysis new\\
Prop.~\ref{prop:RM} & G8-T3c & v008 \texttt{MINIMAL\_RANK.md} & proved, checked $k\le5$ & proved in text\\
Prop.~\ref{prop:small} & (G11-T3 part) & v011 \texttt{MATHEMATICS.md} \S1.7 & proved+checked & finite enumeration, own code\\
Thm.~\ref{thm:W2} & G11-T3 (=G10-C1=G11-C2) & v011 \texttt{MATHEMATICS.md} \S1.7 & proved+checked & text for $p\ge9$; enumeration for $p\le8$\\
Thm.~\ref{thm:rank3} & G11-T4 & v011 \texttt{RANK\_UNIVERSAL.md} & proved & proved in text modulo Prop.~\ref{prop:small}\\
Cor.~\ref{cor:no4} & G11-C3 (G10-C2 false) & v011 \texttt{RANK\_UNIVERSAL.md} & false (G10-C2) & proved in text\\
Thm.~\ref{thm:class35} & G8-T3b, G9-T4 & v008; v009 \S2.2 & proved+checked & re-enumerated exhaustively\\
Thm.~\ref{thm:class6} & G10-T4 & v010 \texttt{BINARY6.md} & proved+checked & source-reported; representatives re-checked\\
Prop.~\ref{prop:geo-n2} & -- & -- & -- & new; proved in text; verified on all 231 classes\\
Cor.~\ref{cor:cert} & -- & -- & -- & new; certificate only, not an algorithm\\
Rem.~\ref{rem:cooperokur} & -- & Cooper--Okur, arXiv:2504.08083, Cor.~3.1 & published theorem (not a release) & translation to cycle codes derived here; consistency-checked\\
Prop.~\ref{prop:heden-n2} & -- & Heden's tail theorem as quoted in Kurz, arXiv:1708.01113 & published theorem (not re-proved) & derived from the cited theorem; not used in any proof\\\bottomrule
\end{longtable}
\endgroup
